\section{Problem Formulation}
\label{sec:problem}

\subsection{Setting}

Let $\agentset = \{\agent_1, \ldots, \agent_{|\agentset|}\}$ be a finite set
of \emph{agents}, $\shapeset = \{\shape_1, \ldots, \shape_{|\shapeset|}\}$ a
finite set of \emph{task shapes}, and
$\tierset = \{\tier_1 \prec \tier_2 \prec \cdots \prec \tier_{|\tierset|}\}$
a totally ordered finite set of \emph{model tiers}, with $\prec$ denoting
``cheaper than''. In \blackrim, $|\agentset| = 17$, $|\shapeset|$ varies
per agent (typically $3$--$5$), and $|\tierset| = 3$ corresponding to haiku,
sonnet, opus. The Cartesian product $\agentset \times \shapeset \times
\tierset$ has $\approx 170$ \emph{cells} (we say $\approx$ because not every
agent has the same number of canonical shapes).

A \emph{dispatch} is a tuple $d = (\agent, \shape, \tier, \cost, q)$ where
$\agent \in \agentset$, $\shape \in \shapeset$, $\tier \in \tierset$, and
$\cost \in \mathbb{R}_{\geq 0}$ is the realised cost in dollars. The
\emph{quality} $q \in \{0, 1\}$ is a Bernoulli indicator of task success
(\cref{sec:eval} discusses how $q$ is observed in production: from explicit
reviewer signals when available, from downstream task closure otherwise).

\subsection{Cost and quality}

Cost is a known deterministic function $\cost : \tierset \times \mathbb{R}^+
\times \mathbb{R}^+ \to \mathbb{R}_{\geq 0}$ of the tier and the realised
input/output token counts, derived from a published rate sheet.

Quality is a stochastic function whose distribution depends on the agent,
the shape, and the tier:
\begin{equation}
  q \mid (\agent, \shape, \tier) \sim \mathrm{Bernoulli}(\theta_{\agent,\shape,\tier}),
  \label{eq:bernoulli}
\end{equation}
with $\theta_{\agent,\shape,\tier} \in [0,1]$ the unknown success
probability for the cell. The advisor's job is to estimate
$\theta_{\agent,\shape,\tier}$ for every cell and select tiers accordingly.

\begin{assumption}[Tier ordering]
\label{assumption:order}
For every $(\agent, \shape)$, the success probabilities across tiers are
weakly increasing in tier strength:
\(
  \theta_{\agent,\shape,\tier_1} \leq \theta_{\agent,\shape,\tier_2} \leq
  \cdots \leq \theta_{\agent,\shape,\tier_{|\tierset|}}.
\)
\end{assumption}

\Cref{assumption:order} is the formal statement of the empirical observation
that more expensive tiers are at least as likely to succeed. It is supported
by Anthropic's published benchmarks \cite{anthropic2024modelcards} and by the
empirical landscape we present in \cref{sec:eval}. The advisor exploits this
ordering to extrapolate: if sonnet's posterior at a cell credibly exceeds
$\qtol$, then opus's also does, so opus is feasible whenever sonnet is.

\subsection{Quality-loss tolerance}

A \emph{quality-loss tolerance} $\qtol \in [0,1]$ is the maximum
allowable degradation in expected quality versus a reference tier
$\tier^* \in \tierset$ (in our deployment, $\tier^*$ is opus). The advisor's
constraint at a cell is
\begin{equation}
  \theta_{\agent,\shape,\tier^*} - \theta_{\agent,\shape,\tier} \leq \qtol,
  \label{eq:tolerance-constraint}
\end{equation}
i.e.\ the candidate tier $\tier$ may not lose more than $\qtol$
absolute quality versus $\tier^*$. We partition $\qtol$ into four
operational tiers, summarised in \cref{tab:tolerance-tiers}.

\begin{table}[h]
\centering
\caption{Quality-loss tolerance tiers used by the advisor. Critical-path
tasks --- those whose downstream blast radius is large --- are pinned to
$\qtol = 0$ (no downgrade ever).}
\label{tab:tolerance-tiers}
\begin{tabular}{lll}
\toprule
\textbf{Class} & \textbf{$\qtol$} & \textbf{Typical workload} \\
\midrule
Critical & $0$              & ADRs, threat modelling, IR orchestration \\
Strict   & $< 0.02$         & Production code review, prompt review \\
Moderate & $< 0.05$         & Multi-file edits, integration tests, judgment routing \\
Lenient  & $< 0.10$         & Lookups, runbook drafting, single-file edits, doc rewrites \\
\bottomrule
\end{tabular}
\end{table}

\subsection{Asymmetric loss}

A successful downgrade saves $\Delta\cost(\tier^*, \tier) =
\cost(\tier^*) - \cost(\tier) > 0$. A failed downgrade incurs $\Delta\cost$
\emph{plus} an opaque downstream cost: the cost of detection, the cost of
repair, the lost trust. Let $\kappa \geq 1$ scale this opaque cost relative
to the saved API spend. The expected loss of choosing $\tier$ rather than
$\tier^*$ at a cell with success probabilities $\theta^*, \theta$ is
\begin{equation}
  \loss(\tier; \theta^*, \theta)
    = -\theta \, \Delta\cost(\tier^*, \tier)
      + (\theta^* - \theta) \cdot \kappa \cdot \Delta\cost(\tier^*, \tier),
  \label{eq:loss}
\end{equation}
where the first term is the expected savings (negative loss) and the second
is the expected regret-of-failure. The decision rule that minimises
\cref{eq:loss} subject to the tolerance constraint
\cref{eq:tolerance-constraint} is
\begin{equation}
  \tier^\dagger
    = \arg\min_{\tier \in \tierset(\qtol)}
      \cost(\tier),
  \label{eq:decision}
\end{equation}
where
\(
  \tierset(\qtol) = \{ \tier : \theta^* - \theta_{\agent,\shape,\tier} \leq \qtol \}
\)
is the feasible set under tolerance $\qtol$. The decision rule is thus
``cheapest tier whose quality stays within tolerance.'' The advisor's job is
to estimate $\tierset(\qtol)$ from data with high credibility.

\subsection{The Bayesian advisor problem}

Each $\theta_{\agent,\shape,\tier}$ is unknown but admits a posterior
distribution $\posterior_{\agent,\shape,\tier}$ given prior $\prior$ and
observation history. The Bayesian decision rule, replacing
\cref{eq:tolerance-constraint} with its posterior-credibility analogue, is
\begin{equation}
  \tier^\dagger
    = \arg\min_{\tier \in \tierset_{\posterior, \qtol, \alpha}}
      \cost(\tier),
  \label{eq:bayesian-decision}
\end{equation}
where
\begin{equation}
  \tierset_{\posterior, \qtol, \alpha}
    = \left\{ \tier : \Pr_{\theta \sim \posterior}\!\big[
        \theta^* - \theta_{\agent,\shape,\tier} \leq \qtol
      \big] \geq 1 - \alpha
    \right\}
  \label{eq:credible-feasible}
\end{equation}
is the credibly-feasible set at confidence level $1 - \alpha$ (we use
$\alpha = 0.05$ throughout, conventionally interpreted as ``95\%
credibility''). \Cref{eq:bayesian-decision} is the decision rule CC-TS
implements; \cref{sec:algorithm} develops the algorithm.

\subsection{What the advisor must produce}

Operationally, the advisor exposes a function
\(
  \mathrm{advise} : (\agent, \shape, \qtol) \mapsto (\tier^\dagger,
  \mathrm{rationale}, \mathrm{cost\_diff})
\)
returning the recommended tier, a human-readable rationale, and the cost
differential versus each tier in $\tierset$. The rationale is non-optional:
operators must be able to audit decisions, and the same surface drives the
\texttt{gt dispatch \texttt{-{}-}advise} CLI integration described in
\cref{sec:system}.
